\section{What the attacker gets: the capability--detectability frontier}
\label{sec:frontier}

An attacker who evades the monitor but learns nothing has not won. To weigh detection against gain, we turned every campaign into exactly the (prompt, answer) pairs it would collect---dataset answers for dataset prompts, the recorded ChatGPT replies for WildChat prompts---fine-tuned a small student on them, and measured GSM8K accuracy \citep{cobbe2021gsm8k}. Appendix~\ref{sec:frontier-full} gives the full protocol and Appendix Figure~\ref{fig:stealth-full} the gain curves.

\paragraph{Two regimes.} Students that already have the skill learn only the answer format. Qwen2.5-0.5B \citep{qwen2024qwen25} gained 9.8 points from 100 math prompts and no more from 900, and padded or natural-looking prompts did as well; Qwen2.5-1.5B on MATH-500 was flat from 1{,}000 to 4{,}000 queries. In this \emph{elicitation} regime, consistent with \citet{gudibande2024false}, stealth is free. A price of stealth can appear only when the student \emph{lacks} the skill. SmolLM2-360M \citep{allal2025smollm2} starts at 3.5\% and gained 7.5, 10.4 and 15.9 points from 1{,}000, 3{,}000 and 10{,}000 MetaMathQA queries---the \emph{extraction} regime. Padding half of every campaign with natural prompts cost 3.6, 2.8 and 3.8 points (48\%, 27\% and 24\% of the gain): the natural prompts taught nothing, so the cost of mimicry is the queries spent on it.

\paragraph{The covert end.} Proposition~\ref{prop:lfc}'s least-detectable campaign---half real WildChat prompts in the math region, half the rest---evaded every detector and gained 2.2 points at 3{,}000 queries, against 10.4: covertness cost 79\% of the gain. To test whether WildChat's math is simply the wrong math, we let the attacker \emph{write} useful prompts in users' voices. An instruction-tuned model rewrote GSM8K-style problems in the style of real WildChat messages, but a classifier still told the rewrites from real prompts with AUC 0.97--0.98, and only 7.9\% passed as user-like. Resampling all rewrites to match real on-task prompts as closely as the classifier could measure gave a campaign that evaded every detector and gained 4.4 points---42\% of the dataset gain, the best any covert campaign achieved---from an effective sample of only 151 distinct rewrites.

\takeaway{The frontier, in one sentence.}{In our testbed, every campaign that evaded every detector also gave up 58\% or more of what a detectable campaign transferred, and each was limited by its supply of prompts that are both useful and indistinguishable from users'.}

This is narrower than it sounds: one 360M student, one benchmark, answers from datasets rather than a live teacher, and two seeds; the covert gains of 1.1--4.4 points are differences of 4--18 of 400 problems. And attackers with a user base do not face the supply limit (Section~\ref{sec:world}).

\paragraph{Capability held at first detection.} What an attacker optimises is the capability it holds when first caught. Table~\ref{tab:timefrontier} combines the gain curves with detection times at 1{,}640 queries a week under two rules applied identically to every detector: the \emph{first weekly alarm} at a 1\% threshold, which is fast but accumulates false alarms (null runs alarmed within 16 weeks a quarter of the time on untouched weeks), and the \emph{anytime-valid} conformal test martingale, which holds the probability of any false alarm, ever, at 1\%.

\begin{table}[t]
\caption{Capability held at first detection: GSM8K gain (points) of the extraction-regime student at the number of queries sent by the median week of first detection, at 1{,}640 queries a week. Each cell: development weeks / untouched weeks. ``none'': no alarm within 16 weeks.}
\label{tab:timefrontier}
\begin{center}
\footnotesize\setlength{\tabcolsep}{3.5pt}
\begin{tabular}{llccc}
\toprule
Detector & Campaign & Weekly detection & First weekly alarm & Anytime-valid \\
\midrule
Two-scale & Dataset & 0.79--0.85 / 0.86 & week 1 (8.8) / week 1 (8.8) & week 3 (12.7) / week 3 (12.7) \\
 & Half padding & 0.21 / 0.00 & week 3 (9.5) / none & week 15 (15.5) / none \\
Hollowed & Dataset & 1.00 / 1.00 & week 1 (8.8) / week 1 (8.8) & not run \\
 & Half padding & 0.44--0.48 / 0.50--0.57 & week 2 (7.9) / week 1 (5.6) & not run \\
Multi-outlier & Dataset & 1.00 / 0.97 & week 1 (8.8) / week 1 (8.8) & not run / week 3 (12.7) \\
 & Half padding & 0.36--0.42 / 0.00 & week 2 (7.9) / none & not run / 25\% of runs \\
Any & Least detectable & $\approx$0.01 / $\le$0.04 & \multicolumn{2}{c}{indistinguishable from false alarms} \\
\bottomrule
\end{tabular}
\end{center}
\end{table}

Against the two-scale and multi-outlier tests, padding on the untouched weeks bought evasion for at least a quarter of a year in most runs at a cost of a quarter to a half of the gain. Against hollowed selection it bought nothing: padded campaigns were caught in the first week. The time to an anytime-valid alarm for hollowed selection is the most important number we did not measure.
